\subtitle{\bf Aula 8 - O Algoritmo LL(1)}
\begin{frame}[t,plain]
    \titlepage
\end{frame}

\begin{frame}{Objetivos}

    \begin{itemize}
        \item
        Apresentar os conceitos de conjuntos first e follow.
        \item
        Apresentar o algoritmo LL(1): primeiro baseado em tabela.
    \end{itemize}
\end{frame}



\begin{frame}{O Nome LL(1)}

    \begin{itemize}

        \item
        \textbf{L} (primeiro): lê a entrada da \textbf{esquerda para a
        direita} (\emph{Left-to-right})
        \item
        \textbf{L} (segundo): produz uma \textbf{derivação mais à esquerda}
        (\emph{Leftmost derivation})
        \item
        \textbf{1}: usa apenas \textbf{1 símbolo de lookahead} para decidir
        qual produção aplicar
    \end{itemize}
\end{frame}

\begin{frame}{A Ideia Central}

    \begin{itemize}

        \item
        Substitui a \textbf{recursão} por um \textbf{algoritmo dirigido por
        tabela}
        \item
        Tabela de parsing \(M[A, a]\): dado não-terminal \(A\) e token \(a\),
        qual produção aplicar?
    \end{itemize}
\end{frame}

\begin{frame}{A Ideia Central}

    \begin{itemize}

        \item
        O analisador mantém uma \textbf{pilha} de símbolos (o que ainda falta
        reconhecer)
        \item
        Cada passo: consulta \(M[\text{topo da pilha}, \text{token atual}]\)
    \end{itemize}
\end{frame}

\begin{frame}{O Algoritmo LL(1)}

    \textbf{Inicialização}: pilha \(\leftarrow [S, \$]\), onde \(S\) é o
    símbolo inicial

    \textbf{Repita}:

    \begin{itemize}

        \item
        Seja \(X\) o topo da pilha e \(a\) o token corrente
        \item
        Se \(X = \$\) e \(a = \$\): \textbf{aceitar}
    \end{itemize}
\end{frame}

\begin{frame}{O Algoritmo LL(1)}

    \begin{itemize}

        \item
        Se \(X\) é terminal \(t\): se \(t = a\), desempilha e avança; senão
        \textbf{erro}
        \item
        Se \(X\) é não-terminal \(A\): consulta \(M[A, a]\)

        \begin{itemize}

            \item
            Se \(M[A, a] = A \to \alpha\): substitui \(A\) na pilha por
            \(\alpha\) (da direita para esquerda)
            \item
            Se \(M[A, a] = vazio\) erro: \textbf{erro sintático}
        \end{itemize}
    \end{itemize}
\end{frame}


\section{\texorpdfstring{Exemplo: Análise de
\texttt{n\ +\ n}}{Exemplo: Análise de n + n}}\label{exemplo-anuxe1lise-de-n-n}

\begin{frame}{Gramática de Expressões (LL(1))}

    \[\begin{array}{lcl}
    E  &\to& T\; E' \\
    E' &\to& +\; T\; E' \mid \lambda \\
    T  &\to& F\; T' \\
    T' &\to& *\; F\; T' \mid \lambda \\
    F  &\to& (\; E\; ) \mid \mathbf{n}
\end{array}\]
\end{frame}

\begin{frame}{Trace Completo da Análise de
    \texttt{n + n}}

    \begin{table}
        \begin{tabular}{l l l l}
            \hline
            Pilha & Entrada & Ação \\
            \hline
            \(E\;\$\) & \(n\;+\;n\;\$\) & \(M[E,n]\): aplicar \(E \to T\;E'\) \\
            \(T\;E'\;\$\) & \(n\;+\;n\;\$\) & \(M[T,n]\): aplicar \(T \to F\;T'\) \\
            \(F\;T'\;E'\;\$\) & \(n\;+\;n\;\$\) & \(M[F,n]\): aplicar \(F \to n\) \\
            \hline
        \end{tabular}
    \end{table}
\end{frame}

\begin{frame}{Trace Completo da Análise de
    \texttt{n + n}}

    \begin{table}
        \begin{tabular}{l l l}
            \hline
            Pilha & Entrada & Ação \\
    \hline
    \(n\;T'\;E'\;\$\) & \(n\;+\;n\;\$\) & Casar: consumir \(n\) \\
    \(T'\;E'\;\$\) & \(+\;n\;\$\) & \(M[T',+]\): aplicar
    \(T' \to \lambda\) \\
    \(E'\;\$\) & \(+\;n\;\$\) & \(M[E',+]\): aplicar \(E' \to +\;T\;E'\) \\
    \hline
\end{tabular}
\end{table}
\end{frame}

\begin{frame}{Trace Completo da Análise de
    \texttt{n + n}}

    \begin{table}
        \begin{tabular}{l l l l}
            \hline
            Pilha & Entrada & Ação \\
            \hline
            \(+\;T\;E'\;\$\) & \(+\;n\;\$\) & Casar: consumir \(+\) \\
            \(T\;E'\;\$\) & \(n\;\$\) & \(M[T,n]\): aplicar \(T \to F\;T'\) \\
            \(F\;T'\;E'\;\$\) & \(n\;\$\) & \(M[F,n]\): aplicar \(F \to n\) \\
            \hline
        \end{tabular}
    \end{table}
\end{frame}

\begin{frame}{Trace Completo da Análise de
    \texttt{n + n}}

    \begin{table}
        \begin{tabular}{l l l l}
            \hline
            Pilha & Entrada & Ação \\
            \hline
            \(n\;T'\;E'\;\$\) & \(n\;\$\) & Casar: consumir \(n\) \\
            \(T'\;E'\;\$\) & \(\$\) & \(M[T',\$]\): aplicar \(T' \to \lambda\) \\
            \(E'\;\$\) & \(\$\) & \(M[E',\$]\): aplicar \(E' \to \lambda\) \\
            \(\$\) & \(\$\) & \textbf{Aceitar} \\
            \hline
        \end{tabular}
    \end{table}
\end{frame}


\section{Construção da Tabela
LL(1)}\label{construuxe7uxe3o-da-tabela-ll1}

\begin{frame}{Conceito: Anulável}

    Um não-terminal \(A\) é \textbf{anulável} se
    \(A \Rightarrow^* \lambda\).

    Uma sequência \(X_1\, X_2\, \cdots\, X_n\) é anulável se \textbf{todos}
    os \(X_i\) são anuláveis.
\end{frame}

\begin{frame}{Conceito: Anulável}

    Casos base:

    \begin{itemize}

        \item
        Todo terminal não é anulável
        \item
        \(A\) anulável se tem produção \(A \to \lambda\)
        \item
        Calculado iterativamente até ponto fixo
    \end{itemize}
\end{frame}

\begin{frame}{Conjunto FIRST}

    \[\mathrm{FIRST}(\alpha) = \{a \in \Sigma \mid \alpha \Rightarrow^* a\,\beta\}\]

    Incluir \(\lambda\) em \(\mathrm{FIRST}(\alpha)\) se \(\alpha\) é
    anulável.
\end{frame}

\begin{frame}{Conjunto FIRST}

    \textbf{Regras para \(\mathrm{FIRST}(X_1\, X_2\, \cdots\, X_n)\):}

    \begin{enumerate}

        \item
        Sequência vazia: \(\mathrm{FIRST}(\lambda) = \{\lambda\}\)
        \item
        \(X_1\) é terminal \(a\): \(\mathrm{FIRST}(X_1\cdots) = \{a\}\)
        \item
        \(X_1\) é não-terminal:

        \begin{itemize}

            \item
            Incluir \(\mathrm{FIRST}(X_1) \setminus \{\lambda\}\)
            \item
            Se \(\lambda \in \mathrm{FIRST}(X_1)\): incluir também
            \(\mathrm{FIRST}(X_2\cdots X_n)\)
        \end{itemize}
    \end{enumerate}
\end{frame}

\begin{frame}{Conjunto FOLLOW}

    \[\mathrm{FOLLOW}(A) = \{a \in \Sigma \cup \{\$\} \mid S \Rightarrow^* \alpha\,A\,a\,\beta\}\]
\end{frame}

\begin{frame}{Conjunto FOLLOW}

    \textbf{Regras:}

    \begin{enumerate}

        \item
        \(\$ \in \mathrm{FOLLOW}(S)\) (símbolo inicial)
        \item
        Para cada produção \(B \to \alpha\,A\,\beta\):

        \begin{itemize}

            \item
            Adicionar \(\mathrm{FIRST}(\beta) \setminus \{\lambda\}\) a
            \(\mathrm{FOLLOW}(A)\)
            \item
            Se \(\lambda \in \mathrm{FIRST}(\beta)\): adicionar
            \(\mathrm{FOLLOW}(B)\) a \(\mathrm{FOLLOW}(A)\)
        \end{itemize}
    \end{enumerate}

    Calculado iterativamente até ponto fixo.
\end{frame}

\begin{frame}{Construção da Tabela}

    Para cada produção \(A \to \alpha\):

    \begin{enumerate}

        \item
        Para cada \(a \in \mathrm{FIRST}(\alpha) \setminus \{\lambda\}\):

        \begin{itemize}

            \item
            Definir \(M[A, a] = A \to \alpha\)
        \end{itemize}
    \end{enumerate}
\end{frame}

\begin{frame}{Construção da Tabela}

    \begin{enumerate}
        \setcounter{enumi}{1}

        \item
        Se \(\lambda \in \mathrm{FIRST}(\alpha)\):

        \begin{itemize}

            \item
            Para cada \(b \in \mathrm{FOLLOW}(A)\):

            \begin{itemize}

                \item
                Definir \(M[A, b] = A \to \alpha\)
            \end{itemize}
        \end{itemize}
    \end{enumerate}

    Entradas não preenchidas indicam \textbf{erro}.
\end{frame}

\begin{frame}{Construção da Tabela}

    \textbf{Conflito} = uma entrada recebe mais de uma produção

    \begin{itemize}

        \item
        gramática não é LL(1)!
    \end{itemize}
\end{frame}


\section{Exemplo Completo: Tabela
LL(1)}\label{exemplo-completo-tabela-ll1}

\begin{frame}{Gramática de Expressões (LL(1))}

    \[\begin{array}{lcl}
    E  &\to& T\; E' \\
    E' &\to& +\; T\; E' \mid \lambda \\
    T  &\to& F\; T' \\
    T' &\to& *\; F\; T' \mid \lambda \\
    F  &\to& (\; E\; ) \mid \mathbf{n}
\end{array}\]
\end{frame}

\begin{frame}{Conjuntos FIRST}

    \begin{table}
        \begin{tabular}{l l l l}
            \hline
            Não-terminal & FIRST \\
            \hline
            \(E\) & \(\{(, n\}\) \\
            \(E'\) & \(\{+, \lambda\}\) \\
            \(T\) & \(\{(, n\}\) \\
            \(T'\) & \(\{*, \lambda\}\) \\
            \(F\) & \(\{(, n\}\) \\
            \hline
        \end{tabular}
    \end{table}
\end{frame}

\begin{frame}{Conjuntos FOLLOW}

    \begin{table}
        \begin{tabular}{l l l l}
            \hline
            Não-terminal & FOLLOW \\
            \hline
            \(E\) & \(\{\$, )\}\) \\
            \(E'\) & \(\{\$, )\}\) \\
            \(T\) & \(\{+, \$, )\}\) \\
            \(T'\) & \(\{+, \$, )\}\) \\
            \(F\) & \(\{*, +, \$, )\}\) \\
            \hline
        \end{tabular}
    \end{table}
\end{frame}

\begin{frame}{Tabela de Parsing LL(1)}

    \begin{table}
        \begin{tabular}{l l l l l l l l}
            \hline
 & \(n\) & \(+\) & \(*\) & \((\) & \()\) & \(\$\) \\
\hline
\(E\) & \(T\;E'\) & --- & --- & \(T\;E'\) & --- & --- \\
\(E'\) & --- & \(+\;T\;E'\) & --- & --- & \(\lambda\) & \(\lambda\) \\
\(T\) & \(F\;T'\) & --- & --- & \(F\;T'\) & --- & --- \\
\(T'\) & --- & \(\lambda\) & \(*\;F\;T'\) & --- & \(\lambda\) &
\(\lambda\) \\
\(F\) & \(n\) & --- & --- & \((\;E\;)\) & --- & --- \\
\hline
\end{tabular}
\end{table}

Cada entrada tem no máximo uma produção --- a gramática é LL(1).
\end{frame}


\section{Conflitos na Tabela LL(1)}\label{conflitos-na-tabela-ll1}

\begin{frame}[fragile]{Conflito FIRST/FIRST}

    Quando duas alternativas de um mesmo não-terminal começam com o mesmo
    terminal:

    \begin{lstlisting}
        S -> if E then S else S
        | if E then S
        | other
    \end{lstlisting}
\end{frame}

\begin{frame}{Conflito FIRST/FIRST}

    \(\mathrm{FIRST}(\mathbf{if}\;E\;\mathbf{then}\;S\;\mathbf{else}\;S) = \{\mathbf{if}\}\)
    \(\mathrm{FIRST}(\mathbf{if}\;E\;\mathbf{then}\;S) = \{\mathbf{if}\}\)

    \(M[S, \mathbf{if}]\) recebe \textbf{duas produções}: \textbf{conflito
    FIRST/FIRST}!
\end{frame}

\begin{frame}[fragile]{Conflito FIRST/FOLLOW}

    Quando uma produção pode derivar \(\lambda\) e os conjuntos FIRST e
    FOLLOW não são disjuntos:

    \begin{lstlisting}
        S -> A a
        A -> a | \lambda{}
    \end{lstlisting}
\end{frame}

\begin{frame}{Conflito FIRST/FOLLOW}

    \begin{itemize}

        \item
        \(\mathrm{FIRST}(A) = \{a, \lambda\}\)
        \item
        \(\mathrm{FOLLOW}(A) = \{a\}\)
    \end{itemize}

    Para \(M[A, a]\): tanto \(A \to a\) (pois \(a \in \mathrm{FIRST}(A)\))
    quanto \(A \to \lambda\) (pois \(a \in \mathrm{FOLLOW}(A)\)):
    \textbf{conflito FIRST/FOLLOW}!
\end{frame}


\section{Implementação em
Haskell}\label{implementauxe7uxe3o-em-haskell}

\begin{frame}[fragile]{Tipos Principais}

    \begin{lstlisting}
        data Symbol = Terminal String | NonTerminal String

        type Production = [Symbol]
        type Grammar    = Map String [Production]

        type ParseTable = Map (String, String) ParseAction

        data ParseAction
        = Derive Production
        | Accept
        | Error
    \end{lstlisting}
\end{frame}

\begin{frame}[fragile]{Cálculo de FIRST por Ponto Fixo}

    \begin{lstlisting}
        computeFirstSets :: Grammar -> FirstSet
        computeFirstSets g = fixedPoint step initial
        where
        initial = Map.fromList [(nt, []) | nt <- nonTerminals g]
        step cur = Map.mapWithKey update cur
        where
        update nt _ = unions $ map (firstForSequence g cur)
        (fromMaybe [] (Map.lookup nt g))
    \end{lstlisting}

    \begin{itemize}

        \item
        \texttt{fixedPoint step initial}: aplica
        \texttt{step} repetidamente até convergência
        \item
        Cada iteração pode expandir os conjuntos FIRST
    \end{itemize}
\end{frame}

\begin{frame}[fragile]{\texttt{firstForSequence}:
    Implementação das Regras}

    \begin{lstlisting}
        firstForSequence :: Grammar -> FirstSet -> [Symbol] -> [String]
        firstForSequence _ _ [] = ["\lambda{}"]
        firstForSequence g fs (sym:syms)
        | isTerminal sym =
        if isLambda sym
        then firstForSequence g fs syms
        else [symbolString sym]
        | otherwise =
        let nt         = symbolString sym
        firstOfSym = fromMaybe [] (Map.lookup nt fs)
        withoutLam = firstOfSym \\ ["\lambda{}"]
        in if "\lambda{}" `elem` firstOfSym
        then withoutLam `union` firstForSequence g fs syms
        else withoutLam
    \end{lstlisting}
\end{frame}

\begin{frame}[fragile]{O Parser LL(1)}

    \begin{lstlisting}
        parse :: Grammar -> ParseTable -> [String] -> ParseResult
        parse grammar table tokens = go initialStack (tokens ++ ["$"]) []
        where
        startSym     = head (nonTerminals grammar)
        initialStack = [NonTerminal startSym, Terminal "$"]
    \end{lstlisting}
\end{frame}

\begin{frame}[fragile]{Parser LL(1)}

    \begin{lstlisting}
        go (Terminal "$" : _) ("$" : _) acc = ParseOk (reverse acc)

        go (Terminal t : stack') (a : inp') acc
        | t == a    = go stack' inp' acc
        | otherwise = ParseError ("expected '" ++ t ++ "'")

        go (NonTerminal nt : stack') inp@(a : _) acc =
        case Map.lookup (nt, a) table of
        Just (Derive prod) ->
        go (filter (not . isLambda) prod ++ stack') inp (prod : acc)
        _ -> ParseError ("no rule for (" ++ nt ++ ", " ++ a ++ ")")
    \end{lstlisting}
\end{frame}


\section{Conclusão}\label{conclusuxe3o}

\begin{frame}{Conclusão}

    \begin{itemize}
        \item
        Apresentamos os conceitos de FIRST, FOLLOW e anulável.
        \item
        Apresentamos a construção da tabela LL(1) e seu algoritmo de parsing.
    \end{itemize}
\end{frame}

\begin{frame}{Conclusão}

    \begin{itemize}

        \item
        Não suporta \textbf{recursão à esquerda} (requer transformação)
        \item
        Exige \textbf{fatoração à esquerda} (prefixos comuns causam conflitos)
        \item
        Algumas gramáticas são inerentemente não-LL(1)
    \end{itemize}
\end{frame}
